\section*{Resumen}

La gestión y consolidación de grandes volúmenes de datos históricos en la Dirección Académica de la Universidad Politécnica de Texcoco ha representado un desafío crítico para la toma de decisiones estratégicas. El manejo disperso de la información mediante hojas de cálculo manuales generaba retrasos operativos, inconsistencias en los registros y dificultaba el seguimiento longitudinal de indicadores clave como la matrícula estudiantil, el aprovechamiento académico, la reprobación, la deserción, la eficiencia terminal y los índices de titulación generacional. Para resolver esta problemática, el presente trabajo de titulación describe el diseño, desarrollo, validación y despliegue de una Plataforma Web de Análisis Estadístico Académico, concebida como una solución centralizada y segura para la ingestión, validación, procesamiento analítico y visualización interactiva de series históricas de datos institucionales.

La plataforma fue desarrollada bajo una arquitectura desacoplada de microservicios contenerizados con Docker y guiada por la metodología ágil Kanban. El backend fue implementado con FastAPI sobre Python 3.12, integrando procesamiento asíncrono en segundo plano con Celery y Redis, y persistencia relacional transaccional en PostgreSQL 16 normalizada mediante SQLAlchemy 2.0 y migraciones con Alembic. El frontend se estructuró como una aplicación de una sola página (SPA) responsiva con React 18, TypeScript, Tailwind CSS y Apache ECharts, gestionando el estado global y asíncrono con Zustand y TanStack Query. El sistema incorpora mecanismos de autenticación y control de acceso basado en roles (RBAC) mediante JSON Web Tokens (JWT) y cifrado criptográfico de contraseñas con Argon2.

Los resultados en el entorno de producción evidencian una reducción sustancial en los tiempos de consolidación analítica, pasando de varios días a respuestas inmediatas en milisegundos con ingestión asíncrona robusta. La suite de pruebas automatizadas alcanzó un 100\% de aprobación en validación de reglas de negocio y cálculo de fórmulas oficiales. Esta plataforma proporciona a las autoridades universitarias una fuente única y confiable de información para la planeación académica, la rendición de cuentas y la mejora continua de la calidad educativa.
